\section{F28379D Device-Realistic Feasibility}
\label{sec:device}

The device-realistic study tests whether the internal ADC architecture can acquire the proposed observations under register-realizable timing and an explicit high-common-mode analog front end. Four synchronized 16-bit differential ADC modules use a 320-ns aperture and sustain 1.093~MS/s per module. A one-edge schedule places three samples on each side of the reconstructed edge within $T_{w,R}=2.2$~\si{\micro\second} after a 0.5-\si{\micro\second} guard. Direct ADC connection is infeasible over the simulated common-mode range, so feasibility presumes the specified differential AFE, an edge-band common-mode rejection ratio (CMRR) bound, an external reference and clock, and a bounded calibration residual.

Seven representative conditions with 200 Monte Carlo seeds each yield worst-case p95 capacitance and ESR errors of $1.6425\%$ and $2.6109\%$; every condition meets its predeclared threshold. Table~\ref{tab:f28379d} summarizes the evidence for the feasibility claim. Register timing, aperture margin, direct-memory-access traffic, utilization, AFE component assumptions, common-mode extrema, and coverage ranges are retained in Supplementary Note S1.

\begin{table}[t]
\caption{F28379D Device-Realistic Summary}
\label{tab:f28379d}
\centering
\footnotesize
\begin{tabular}{ll}
\toprule
Item & Specified value \\
\midrule
ADC architecture & four modules, 16-bit differential \\
Aperture/rate & 320 ns / 1.093 MS/s per module \\
Edge acquisition & 0.5-\si{\micro\second} guard; $T_{w,R}=2.2$~\si{\micro\second}; 3 points/side \\
AFE boundary & differential path; edge-band CMRR $\ge70$~dB \\
Health reporting & 1024 cycles (20.48 ms) \\
Worst p95 C/ESR & 1.6425\% / 2.6109\% \\
\bottomrule
\end{tabular}
\end{table}

The $1.36$-\si{\micro\second} \TSRLS{} value is an estimator-arithmetic estimate, not platform WCET. Firmware compilation, linker placement, measured execution time, and bench AFE verification remain open; the result supports a constrained simulation architecture rather than hardware validation.


